%% kapa: exam=1 round=19 year=2008 period=1 form=A subject=appraisal_law subject_ko=부동산관계법규 seq=12 number=52
%% source: Q-Net 감정평가사 기출문제 (cst00309) artlSeq=1001962 file=extracted/1008/1001962/2008년 제19회 감정평가사 1차시험 기출문제/1교시.hwp
%% conversion: kordoc 4.24.51 -> tools/gen1.py (auto); review=proofread-sonnet; flags=none
\begin{question}{52}
국토의 계획 및 이용에 관한 법령에 의한 기간 중 가장 짧은 것은? (단, 특별한 사유가 있는 경우 등 예외적인 경우와 심의 또는 협의기간은 고려하지 않음)
\begin{choices}
  \item 개발행위허가신청에 대한 처분기간
  \item 토지거래허가신청의 처분에 대한 이의신청제기기간
  \item 도지사의 도시관리계획 지형도면에 대한 승인기간
  \item 토지거래허가신청의 불허가처분에 대한 매수청구기간
  \item 도시관리계획의 결정을 위한 협의요청을 받은 기관의 의견제시기간
\end{choices}
\end{question}
\answer{1}
